\chapter{Prediction with Gradient Boosting}
\label{ch:xgb}

In the European regime the target of Chapter~\ref{ch:success} --- winsorized EBITDA
arc-growth over a two-year horizon --- is observable for essentially every firm-year, so
the first modelling layer is plain supervised learning. This chapter fixes the learning
problem, the feature catalogue, the validation design, and --- the part practitioners most
often skip --- the uncertainty quantification. The guiding principle throughout is that
firm-year panels are \emph{not} i.i.d.\ samples: rows recur within firms and cluster
within macro eras, and every design choice below (splits, calibration sets, bootstrap
units) follows from taking that dependence seriously. What the fitted model \emph{means}
--- association, not effect --- is deferred to Chapter~\ref{ch:causal}.

\section{The learning problem}
\label{sec:xgb-problem}

A row is a firm-year $(i,y)$. The label is
$Y_{i,y} = w_c\!\bigl(G^{(h)}_{i,y}\bigr)$, the arc-growth of EBITDA between statement
years $y$ and $y+h$ with $h=2$, winsorized at level $c$ exactly as defined in
Chapter~\ref{ch:success}. The feature vector $X_{i,y} \in \mathbb{R}^p$ is assembled
\emph{point-in-time}: it may use only filings published and macro prints released by the
information date $t_{\mathrm{info}}(i,y)$ of the row, respecting the publication-lag
discipline of Chapter~\ref{ch:data}. We estimate two objects:
\begin{equation}
m(x) = \E[Y \mid X = x],
\qquad
q_\tau(x) = \inf\{a : \Prob(Y \le a \mid X = x) \ge \tau\},
\label{eq:xgb-targets}
\end{equation}
the conditional mean (point screening) and a small grid of conditional quantiles. We
use $\tau \in \{0.05, 0.10, 0.50, 0.90, 0.95\}$. The quantiles are not decoration: the
outcome is heavy-tailed even after winsorization, upside screening acts on $q_{0.9}$
rather than the mean, downside risk control acts on $q_{0.1}$, and the pair
$(q_{0.05}, q_{0.95})$ is the raw material for the conformal intervals of
Section~\ref{sec:xgb-uq}.

Two dependence structures dominate. \emph{Within-firm}: rows $(i,y)$ and $(i,y')$ share
persistent firm effects (management quality, niche economics), so errors are serially
correlated within $i$. \emph{Within-era}: rows of a given year $y$ share the macro shock,
so errors are cross-sectionally correlated within $y$. The first dictates firm-grouped
resampling; the second dictates era-blocked validation. Neither is optional.

\section{Feature catalogue}
\label{sec:xgb-features}

Table~\ref{tab:xgb-features} lists the feature families. All are computable from the
fundamentals panel $F_{i,y}$, the deals table where coverage exists, and the macro feed;
all obey the as-of rule.

\begin{table}[htbp]
\centering
\caption{Feature catalogue for the EU prediction layer. All features are as-of
$t_{\mathrm{info}}(i,y)$; macro features use publication-lagged vintages
(Chapter~\ref{ch:data}).}
\label{tab:xgb-features}
\small
\begin{tabular}{@{}lll@{}}
\toprule
Family & Examples & Construction \\
\midrule
Lagged levels & EBITDA, revenue, payroll, assets & $\log(1+\cdot)$ where positive; \\
 & at $y, y{-}1, y{-}2$ & raw for signed EBITDA \\
Growth rates & 1--3y arc-growth of revenue, & same arc transform as the \\
 & EBITDA, payroll & target, unwinsorized \\
Ratios & EBITDA margin, leverage, & clipped at 1st/99th pctile \\
 & asset turnover, payroll share & per sector-year \\
Peer position & within-sector-year rank and & computed per $(s,y)$ cell, \\
 & z-score of margin, growth, size & rank in $[0,1]$ \\
Sector-year aggregates & median growth, IQR, firm count & cell-level, merged back \\
Firm statics & age at $y$, size bucket, country & age in years, log assets bucket \\
Macro as-of & policy rate level and 12m change, & vintage at $t_{\mathrm{info}}$; \\
 & CPI YoY, unemployment level/change & level \emph{and} change \\
Deal history & raised-to-date, round count, & zero-filled with indicator \\
 & months since last round & where firm never in deals \\
\bottomrule
\end{tabular}
\end{table}

\paragraph{Missingness.}
Fundamentals gaps are informative: small companies file abbreviated accounts, distressed
companies file late, and the propensity to be missing depends on the (unobserved) value
--- textbook MNAR. We therefore do \emph{not} mean-impute: imputation overwrites exactly
the signal carried by absence and biases every downstream split. Instead we rely on the
native sparsity handling of gradient boosting, in which each split learns a default
direction for missing values from the data \cite{chen2016}, and we add explicit
indicators (\texttt{ebitda\_missing\_lag1}, \texttt{never\_in\_deals}) for gaps whose
economic meaning we can articulate. Trees are then free to interact the indicator with
size and sector --- ``missing EBITDA at small size'' is a different animal from ``missing
EBITDA at large size'', and the model should be allowed to say so.

\section{Why boosted trees here --- and what they cannot do}
\label{sec:xgb-why}

The honest case for gradient-boosted trees on this problem is fourfold: (i) the signal is
riddled with interactions (margin $\times$ sector, growth $\times$ rate regime) that we
cannot enumerate a priori and trees discover automatically; (ii) the features are mixed
continuous/ordinal/categorical with wildly different scales, and trees are invariant to
monotone transforms of each feature, so we spend no effort on marginal calibration;
(iii) native missingness handling, as above; (iv) strong, well-understood regularization
(depth, shrinkage, subsampling, $\ell_2$ on leaf weights) \cite{chen2016}.

The honest case against has two entries, and both matter here. First, \emph{no
extrapolation}: a tree ensemble is piecewise constant, so its prediction outside the
convex hull of the training features is the value at the nearest boundary cell. If every
training era has falling policy rates, the fitted rate response is a flat extension ---
the model literally cannot form an opinion about a rising-rate regime it has never seen.
Second, \emph{unstable attribution}: correlated features split importance
unpredictably across refits, so feature rankings are not portable conclusions.

We mitigate rather than solve. Era-aware validation (Section~\ref{sec:xgb-validation})
at least measures degradation across regimes instead of hiding it. Monotone constraints
encode signs we are willing to defend economically --- e.g.\ growth non-increasing in the
12-month rate \emph{change}, downside quantile non-increasing in leverage --- which
reduces variance in thin regions and prevents sign flips there.

\begin{remark}[Constraints are priors, not facts]
\label{rem:xgb-mono}
A monotone constraint changes the model class, not the world. Impose one only where the
sign survives scrutiny (rates \emph{level} is ambiguous: it proxies both financing cost
and the demand cycle; the \emph{change} is the cleaner candidate). A wrong constraint is
worse than none, because it converts local noise into global bias silently.
\end{remark}

For the conditional-mean model, heavy tails argue against squared error even after
winsorization. We use the pseudo-Huber objective
\begin{equation}
\ell_\delta(r) = \delta^2\left(\sqrt{1 + (r/\delta)^2} - 1\right),
\label{eq:xgb-phuber}
\end{equation}
which is quadratic for $|r| \ll \delta$ and linear for $|r| \gg \delta$, with $\delta$
set to roughly one cross-sectional standard deviation of $Y$. Note the fitted functional
then sits between conditional mean and median; we report which loss produced which model
and never mix them in evaluation.

\section{Validation design}
\label{sec:xgb-validation}

\paragraph{Never random K-fold.}
Random row-level K-fold fails three ways on this panel. (i)~\emph{Firm leakage}: rows of
the same firm land on both sides of the split, and the persistent firm effect makes
validation loss optimistic --- the model is partly memorizing firms. (ii)~\emph{Target
overlap}: with $h=2$, the labels of years $y$ and $y+1$ share a statement year, so
adjacent rows carry overlapping information about each other's targets.
(iii)~\emph{Era leakage}: the model trains on the future macro regime it will be asked to
predict, which no deployed model enjoys. All three inflate measured skill precisely where
the deployment risk is largest.

\paragraph{Rolling-origin splits.}
We order eras (label years) $y_1 < \dots < y_E$ by information date and build folds
$k = 1, \dots, K$: train on eras up to $y_{e_k}$, \emph{embargo} the next $h$ eras (their
labels overlap the training window), validate on the era after the embargo. Every fold
mimics deployment: fit on the past, predict a disjoint future. The same firm may appear
in train (early years) and validation (later years) --- that is legitimate and realistic;
the leakage to prevent is label overlap, handled by the embargo, not firm recurrence
across time.

\paragraph{Hyperparameter protocol.}
Random search (100--200 draws) over the ranges in Table~\ref{tab:xgb-hparams}; each draw
is scored by its \emph{average validation loss across all rolling folds}, with the number
of boosting rounds chosen by early stopping on the \emph{most recent} fold only ---
recency-weighting the one knob that matters most. The final model refits on all data with
the tuned configuration and the early-stopped round count scaled by the ratio of full to
training rows. Tuning on an average across eras deliberately trades a little peak skill
for regime stability.

\begin{table}[htbp]
\centering
\caption{Search ranges. Depth is kept shallow: firm-year panels reward many weak
interactions over few deep ones.}
\label{tab:xgb-hparams}
\begin{tabular}{@{}lll@{}}
\toprule
Parameter & Range & Note \\
\midrule
max depth & 3--6 & log-uniform integer \\
learning rate & 0.01--0.1 & log-uniform \\
min child weight & 5--100 & guards thin sector cells \\
row subsample & 0.6--0.9 & per boosting round \\
column subsample & 0.5--0.9 & decorrelates trees \\
$\ell_2$ leaf penalty & 0.5--20 & log-uniform \\
rounds & early stopping & patience 50, most recent fold \\
\bottomrule
\end{tabular}
\end{table}

\section{Uncertainty quantification}
\label{sec:xgb-uq}

Three tools, three targets: per-firm heteroskedastic intervals (quantile regression),
finite-sample coverage repair (conformal), and model-level uncertainty of scalar
summaries (cluster bootstrap). Table~\ref{tab:xgb-uqtools} states when each is the right
instrument.

\subsection{Quantile regression with the pinball loss}

For $\tau \in (0,1)$ the pinball (check) loss is
\begin{equation}
\rho_\tau(u) = u\,\bigl(\tau - \ind{u < 0}\bigr)
= \tau\, u_+ + (1-\tau)\, u_-,
\label{eq:xgb-pinball}
\end{equation}
with $u_+ = \max(u,0)$, $u_- = \max(-u,0)$. Its population minimizer is exactly the
quantile we want:

\begin{proposition}[The pinball minimizer is the $\tau$-quantile]
\label{prop:xgb-pinball}
Let $Y$ have distribution function $F$ and finite first moment, and let
$\varphi(a) = \E[\rho_\tau(Y - a)]$. Then any $a^*$ with $F(a^{*-}) \le \tau \le F(a^*)$
minimizes $\varphi$; if $F$ is continuous and strictly increasing,
$a^* = F^{-1}(\tau)$ uniquely.
\end{proposition}

\begin{proof}
$\varphi$ is convex, and for $a$ a continuity point of $F$,
$\varphi'(a) = -\tau \Prob(Y > a) + (1-\tau)\Prob(Y < a) = F(a) - \tau$.
A convex function is minimized where its (sub)derivative changes sign, i.e.\ where
$F(a^{-}) - \tau \le 0 \le F(a) - \tau$.
\end{proof}

We fit one boosted model per $\tau$ with objective \eqref{eq:xgb-pinball} applied
conditionally. Separately fitted quantiles can cross in finite samples; we repair by
sorting the predicted grid per row (monotone rearrangement), which never worsens pinball
loss.

\subsection{Conformalized quantile regression}

Quantile models are calibrated only insofar as the model class and sample size allow;
conformal prediction repairs marginal coverage with finite-sample guarantees
\cite{romano2019}.

\paragraph{Algorithm (split CQR).}
\begin{enumerate}
\item Split the available rows into a proper training set $I_1$ and a calibration set
$I_2$ of size $n_2$, with the split at the firm level and $I_2$ chosen as a
\emph{block}: the most recent complete era(s) before the deployment date.
\item On $I_1$, fit lower and upper quantile models $\hat q_{\alpha/2}$,
$\hat q_{1-\alpha/2}$ as above.
\item For each $j \in I_2$ compute the conformity score
$E_j = \max\bigl\{\hat q_{\alpha/2}(X_j) - Y_j,\; Y_j - \hat q_{1-\alpha/2}(X_j)\bigr\}$.
\item Let $Q_{1-\alpha}$ be the $\lceil (n_2+1)(1-\alpha) \rceil$-th smallest of
$\{E_j\}_{j \in I_2}$.
\item Report the interval
$C(x) = \bigl[\hat q_{\alpha/2}(x) - Q_{1-\alpha},\;
\hat q_{1-\alpha/2}(x) + Q_{1-\alpha}\bigr]$.
\end{enumerate}

\begin{theorem}[CQR marginal coverage \cite{romano2019}]
\label{thm:xgb-cqr}
If the calibration pairs $\{(X_j, Y_j)\}_{j \in I_2}$ and the test pair $(X, Y)$ are
exchangeable, then $\Prob\{Y \in C(X)\} \ge 1 - \alpha$; if the scores are almost surely
distinct, also $\Prob\{Y \in C(X)\} \le 1 - \alpha + 1/(n_2 + 1)$.
\end{theorem}

\begin{warning}[Temporal data are not exchangeable]
\label{warn:xgb-exch}
Theorem~\ref{thm:xgb-cqr} assumes the test row is exchangeable with the calibration
rows. A macro-era shift between calibration block and deployment breaks this, and no
finite-sample guarantee survives. Practice: calibrate on the most recent contiguous
block (as above), \emph{treat the nominal level as approximate}, measure empirical
coverage per era on every historical fold, and recalibrate at every model vintage. If
per-era coverage drifts systematically with the macro state, widen intervals by the
worst observed era's deficit rather than trusting the nominal $\alpha$.
\end{warning}

\subsection{Cluster bootstrap over firms}

Quantile and conformal intervals describe \emph{per-row} outcome uncertainty. A
different question --- how uncertain is a scalar summary of the fitted model, such as an
era MAE, a decile-spread statistic, or the skill gap over a benchmark --- requires
resampling the fitting process itself \cite{efron1994}.

\paragraph{Algorithm (firm-level percentile bootstrap).}
\begin{enumerate}
\item Let $\mathcal{F}$ be the set of firms. For $b = 1, \dots, B$ (we use
$B \in [200, 500]$): draw $|\mathcal{F}|$ firms from $\mathcal{F}$ with replacement and
assemble the replicate panel from \emph{all rows} of each drawn firm (duplicates kept).
\item Re-run the frozen pipeline on the replicate: feature construction, the tuned
hyperparameter configuration (not re-searched), fit, and compute the scalar summary
$\hat\theta_b$ on the replicate's validation eras.
\item Report the percentile interval
$\bigl[\hat\theta_{(\lceil \alpha B/2 \rceil)},\;
\hat\theta_{(\lceil (1-\alpha/2) B \rceil)}\bigr]$.
\end{enumerate}
Resampling firms, not rows, is the load-bearing choice: within-firm correlation means
the effective sample size is closer to the number of firms than the number of rows, and
a row-level bootstrap understates variance accordingly.

\begin{table}[htbp]
\centering
\caption{Choosing the uncertainty tool.}
\label{tab:xgb-uqtools}
\begin{tabular}{@{}llll@{}}
\toprule
Tool & Quantifies & Guarantee & Cost \\
\midrule
Quantile regression & per-row outcome spread & none per se & 1 fit per $\tau$ \\
Split CQR & per-row interval coverage & finite-sample marginal, & 1 calibration \\
 & & approximate under drift & pass \\
Firm bootstrap & model-level scalars & asymptotic, within & $B$ refits \\
 & (MAE, skill gaps) & the observed design & \\
\bottomrule
\end{tabular}
\end{table}

\section{Evaluation protocol}
\label{sec:xgb-eval}

We report MAE and $R^2$ \emph{by era} and by horizon, never pooled alone: a pooled MAE
averages over exactly the regime heterogeneity the model is weakest on. The median model
(pinball at $\tau = 0.5$) is evaluated with MAE --- matched loss --- and the pseudo-Huber
mean model with both MAE and RMSE. Quantile calibration is checked per era with
reliability diagrams: for each $\tau$, the empirical fraction of outcomes below
$\hat q_\tau(X)$ plotted against $\tau$; a well-calibrated model tracks the diagonal, a
regime-sensitive one shows era-dependent slope. Subgroup stability: per-cell MAE by
sector and size tercile relative to pooled MAE, flagging cells where a benchmark of
Section~\ref{sec:xgb-bench} beats the boosted model --- such cells are candidates for
exclusion from deployment, not for silent averaging.

\begin{warning}[Attributions are associational]
\label{warn:xgb-shap}
SHAP-style attributions decompose the \emph{fitted} $\hat m$, i.e.\ they redistribute
predicted differences among correlated inputs of an associational object. A large
attribution on \texttt{months\_since\_last\_round} does not mean financing drives
growth; it means firms that recently raised look different, which is precisely the
selection effect Chapter~\ref{ch:causal} is built to separate from the causal one. Use
attributions for debugging (leakage detection, sanity of signs) and never in an
investment memo as an effect size.
\end{warning}

\section{Benchmarks and when not to use gradient boosting}
\label{sec:xgb-bench}

Every fold carries three benchmarks: \emph{persistence} (predict the firm's last
observed growth), \emph{sector-year mean} (predict the cell average), and \emph{ridge}
regression on the identical feature matrix (with the missingness indicators, since ridge
has no native sparsity handling). The shipping rule is strict: the boosted model is
deployed only where its out-of-era skill gap over ridge exceeds the half-width of the
firm-bootstrap interval on that gap. A gap inside the interval is noise, and ridge is
cheaper, more stable, and partially extrapolates.

\begin{table}[htbp]
\centering
\caption{When gradient boosting is the wrong tool.}
\label{tab:xgb-decision}
\begin{tabular}{@{}lll@{}}
\toprule
Situation & Symptom & Use instead \\
\midrule
Tiny sample & $\lesssim 5$k firm-years & ridge on curated features \\
Regime extrapolation & deployment macro outside & linear/constrained macro \\
 & training support & term $+$ trees on the rest \\
Regulatory interpretability & auditable decision rule required & scorecard or shallow \\
 & & monotone model \\
Skill gap within noise & bootstrap CI covers zero & ridge; revisit features \\
\bottomrule
\end{tabular}
\end{table}

The model of this chapter answers ``which firms will grow?'' conditional on everything
observable. It is silent on ``what happens if a firm raises?'' --- the fitted dependence
on deal-history features mixes selection with effect. Chapter~\ref{ch:causal} takes up
that question with the same data and an explicitly causal contract.
